% Figure: the magnetocaloric refrigeration cycle on the entropy-temperature plane
% of a magnetic solid. Two curves of total entropy against temperature, one at
% low applied field and one at high applied field; the high-field curve lies
% below because aligning the spins removes magnetic entropy. The cycle is an
% analog of the Brayton (or Ericsson with regeneration) cycle: adiabatic
% magnetisation heats, isofield heat rejection, adiabatic demagnetisation cools,
% isofield heat absorption. Drawn as a rectangle-like loop between the two curves.
\begin{figure}[htbp]
\centering
\begin{tikzpicture}
\begin{axis}[
  width=12cm, height=8.4cm,
  xlabel={temperature $T$ (K)},
  ylabel={entropy $s$ (arb.)},
  xmin=270, xmax=320, ymin=0.2, ymax=1.05,
  axis lines=left,
  legend pos=south east, legend cell align=left,
  legend style={font=\scriptsize},
  tick label style={font=\scriptsize}, label style={font=\small},
  clip=false,
]
% zero-field entropy curve: s rises with T (schematic smooth rise)
\addplot[engblue, very thick, samples=80, domain=270:320]
  {0.30 + 0.55/(1 + exp(-(x-295)/6))};
\addlegendentry{low field $H_0$}
% high-field entropy curve: shifted down (spins aligned, less magnetic entropy)
\addplot[engred, very thick, samples=80, domain=270:320]
  {0.20 + 0.55/(1 + exp(-(x-303)/6))};
\addlegendentry{high field $H_1$}
% cycle loop: adiabatic magnetisation (vertical up in T at const s) from A to B,
% isofield reject B to C (down the high-field curve), adiabatic demag C to D,
% isofield absorb D to A.
% A: on low-field curve at T=290 : s = 0.30 + 0.55/(1+exp((295-290)/6)) approx
% pick simple loop points near the curves.
% A (287, 0.44) low field ; adiab magnetise to B (300, 0.44) at const s on high field
\draw[engblack, very thick, -{Latex}] (axis cs:287,0.44) -- (axis cs:300,0.44);
\node[font=\scriptsize, anchor=south] at (axis cs:293.5,0.445) {adiab. magnetise (warms)};
% B (300,0.44) reject heat along high-field curve down to C (293,0.31)
\draw[engblack, very thick, -{Latex}] (axis cs:300,0.44) -- (axis cs:293,0.31);
\node[font=\scriptsize, anchor=west] at (axis cs:300,0.38) {reject $q_h$};
% C (293,0.31) adiab demagnetise at const s to D (280,0.31) on low-field curve
\draw[engblack, very thick, -{Latex}] (axis cs:293,0.31) -- (axis cs:280,0.31);
\node[font=\scriptsize, anchor=north] at (axis cs:286.5,0.305) {adiab. demag. (cools)};
% D (280,0.31) absorb heat along low-field curve up to A (287,0.44)
\draw[engblack, very thick, -{Latex}] (axis cs:280,0.31) -- (axis cs:287,0.44);
\node[font=\scriptsize, anchor=east] at (axis cs:280,0.38) {absorb $q_c$};
\addplot[engblack, only marks, mark=*, mark size=1.2pt] coordinates
  {(287,0.44) (300,0.44) (293,0.31) (280,0.31)};
\node[font=\scriptsize, anchor=east] at (axis cs:287,0.44) {A};
\node[font=\scriptsize, anchor=west] at (axis cs:300,0.44) {B};
\node[font=\scriptsize, anchor=west] at (axis cs:293,0.31) {C};
\node[font=\scriptsize, anchor=east] at (axis cs:280,0.31) {D};
\end{axis}
\end{tikzpicture}
\caption[Magnetocaloric refrigeration cycle on the entropy-temperature plane]{The
magnetocaloric cycle plotted on the entropy-temperature plane of a magnetic
solid. Two curves give the material's total entropy against temperature, one at
low applied magnetic field and one at high field; the high-field curve lies lower
because aligning the magnetic spins removes magnetic entropy. Applying the field
adiabatically (A to B) holds entropy fixed, so the temperature rises, the
magnetocaloric heating; the solid then rejects heat to the warm side along the
high-field curve (B to C); removing the field adiabatically (C to D) drops the
temperature, the magnetocaloric cooling; and the solid absorbs heat from the cold
load along the low-field curve (D to A). The cycle is the magnetic analog of the
gas refrigeration cycles of this chapter, with the applied field playing the role
of pressure, and near a material's magnetic ordering temperature the entropy
change with field is large enough to make room-temperature magnetic refrigeration
a serious candidate for replacing the vapour-compression cycle without any
volatile refrigerant.}
\label{fig:vol04ch03:magnetocaloric-st}
\end{figure}
